Exotic decays of the Higgs boson are well motivated by a variety of models beyond the Standard Model (BSM) \cite{Curtin_2014}, including supersymmetry, Dark Sector or extended Higgs sectors, and provide a unique opportunity to probe new physics.  

This paper presents a new strategy to search for exotic Higgs boson decays to a final state with a photon ($\gamma$) and missing transverse momentum (\ETmiss), probing an unexplored region of the  phase space with low kinematic thresholds. 
The considered signature can arise %as the rare $H\to Z(\nu\nu)\gamma$ decay predicted by the SM, but also 
in several BSM models, such as models with very low-scale supersymmetry breaking where the Higgs boson can decay into a gravitino ($\tilde{G}$) and a bino ($\tilde{B}$) that subsequently decays via $\tilde{B}\to\gamma\tilde{G}$ \cite{Petersson_2012}, in the Peccei-Quinn limit of Next to Minimal Supersymmetric Models (NMSSM) \cite{Suematsu_1998}, or in the decay of the Higgs boson into a pair of right-handed neutrinos (N) subsequently decaying via $N\to\nu\nu$ and $N\to\gamma N$ \cite{graesser2007}. 
In addition, this semi-visible decay of the Higgs boson provides an interesting production channel for the dark photon ($\gammad$), which arises as the gauge boson of an additional \( U(1)_D \) gauge group in Dark Sector models. Strong motivations for the search of this particle arise from its potential contribution in solving the small-scale structure formation problems \cite{Fan:2013tia} \cite{Heikinheimo:2015kra}, as well as being a possible source for the Sommerfeld enhancement of Dark Matter (DM) annihilation required to explain the PAMELA-Fermi-MS2 positron anomaly in terms of DM \cite{Arkani-Hamed:2008hhe}, or making asymmetric DM scenarios viable \cite{Zurek:2013wia}. 

While \gammad\  searches typically target the massive \gammad\ \cite{Beacham_2019,harris2022snowmasswhitepapernew,Deliyergiyev_2016}, which can couple directly to the Standard Model (SM) via kinetic mixing with the ordinary photon \cite{Fabbrichesi_2021}, this paper focuses on the massless \gammad\ scenario. %\textcolor{red}{but the results remain valid also for a dark photon of mass up to around 10 GeV, as previous studies have shown that the kinematics are not sensitive to the mass of the dark photon up to these values ~\cite{EXOT-2021-17}}. 
In this case, the kinetic mixing between the \gammad\ and the photon can be tuned away through a rotation of the gauge fields, and the \gammad\ can interact with ordinary matter only through operators of dimension higher
than four \cite{Fabbrichesi_2021}. The Higgs boson can decay through a loop interaction involving additional (pseudo-)scalar BSM messenger fields, which couple both to the SM and the Dark Sector~\cite{Gabrielli:2014oya, Biswas:2016jsh}, as shown in 
%Fig.~\ref{fig:Feynman-ggF}
Fig.~\ref{fig:Feynman-all-Hyyd}. Thanks to the non-decoupling properties of the Higgs boson, and depending on the parameters of the model, a branching ratio of $H\to\gamma\gammad$ (\br) of few percent can be allowed ~\cite{Beauchesne2023a, Beauchesne2023b}. 

% \begin{figure}[htbp]
%     \centering
%     \includegraphics[width=0.4\textwidth]{figures/Feynman_ggF.png}
%     \caption{Feynman diagram illustrating the gluon-gluon fusion (ggF) process.}
%     \label{fig:Feynman-ggF}
% \end{figure}

\begin{figure}[htbp]
    \centering
    \subfloat[]{\includegraphics[width=0.4\textwidth]{figures/feynman_ggHyyd.pdf}\label{ggh-diagram}}

    \subfloat[]{\includegraphics[width=0.4\textwidth]{figures/feynman_VBFHyyd.pdf}\label{vbf-diagram}}
    \hspace{.1\textwidth}
    \subfloat[]{\includegraphics[width=0.4\textwidth]{figures/feynman_VHyyd.pdf}\label{vh-diagram}}
    \caption{Feynman diagrams illustrating the Higgs boson production via gluon-gluon fusion process (a), vector-boson fusion (b), and Higgsstrahlung (c), all followed by the decay into a photon and a dark photon ($H\to\gamma\gammad$). The hatched circle in the $H\gamma\gammad$ vertex represents the interaction through messenger fields.}
    \label{fig:Feynman-all-Hyyd}
\end{figure}

Searches for the semi-visible decay $H\to\gamma\gammad$  have been conducted extensively at the LHC by the ATLAS and CMS experiments, in the vector boson fusion (VBF) Higgs production channel and in the associated production with $Z$ boson ($ZH$)~\cite{HDBS-2019-13, EXOT-2021-17, EXOT-2018-63, CMS-EXO-19-007,EXOT-2014-06,CMS-EXO-12-047,CMS-HIG-14-025,CMS-EXO-20-005}. The combination of ATLAS results in the $ZH$ and VBF channels provided the best limit so far, excluding a \br> 1.3\% (1.5\% expected) at 95\% CL with 139 fb\(^{-1}\) of data~\cite{ATLAS:2024cju} collected during the LHC Run-2 at $\sqrt{s}=13$\ TeV. 
%CMS probed this decay using Higgs events produced in association with a \( Z \) boson (\( ZH \), \( Z \to \ell^+\ell^- \)) with 137 fb\(^{-1}\), setting a 95\% CL upper limit of 4.6\% (3.6\% expected), while ATLAS achieved a limit of 2.8\% (2.3\% expected) with 139 fb\(^{-1}\)~\cite{HDBS-2019-13} in the same final state. In vector-boson fusion (VBF) production with 130 fb\(^{-1}\), CMS achieved a 3.5\% (2.8\% expected) limit. ATLAS set a more stringent 1.8\% (1.7\% expected) limit in the VBF channel with 139 fb\(^{-1}\)~\cite{EXOT-2021-17}. The combination of ATLAS results for ZH and VBF leads to a combined limit of 1.3\% (1.5\% expected) with 139 fb\(^{-1}\) of data~\cite{ATLAS:2024cju}. 
Though these two production modes provide clean final states, the gluon-gluon fusion one (ggF), leading to the $\gamma+\ETmiss$ final state, is one order of magnitude larger in cross-section than the VBF channel. This final state has been explored by CMS ~\cite{CMS-HIG-14-025} and ATLAS in the context of other exotic models~\cite{EXOT-2018-63} or focusing on the decay into $\gamma\gammad$ of BSM Higgs bosons with mass higher than 400 GeV~\cite{ATL-PHYS-PUB-2023-003}. However, the SM $H \to \gamma \gamma_d$ decay has not been specifically targeted, as the low thresholds on the photon transverse momentum (\pTy) and \ETmiss\  required for this signature constitute a challenge for trigger efficiency. Developments in the trigger between Run-2 and Run-3 led to the implementation of a composite trigger that combines selections on \ETmiss, \pTy, and the transverse mass between the $\gamma$ and the \ETmiss\ ($\mT=\sqrt{2\pTy\ETmiss(1-\cos(\Delta\phi))}$) allowing lower thresholds and increasing the acceptance of ggF events by approximately a factor of two, while keeping the trigger rate under control. By exploiting this new trigger, the analysis presented in this paper is optimised for the $H\to\gamma\gammad$ decay in the ggF production channel for the first time in ATLAS, retaining sensitivity also to the VBF, $ZH$ and $WH$ channels. 
